\documentclass[12pt]{article}
\usepackage[utf8]{inputenc}
\usepackage[T1]{fontenc}
\usepackage{hyperref}
\title{Polish Sign Language Language Resources for K-12 Curriculum Development}
\author{Compiled by Claude Sonnet 5.5 (claude-sonnet-5-5) for Andr\'as Kornai}
\date{}
\begin{document}
\maketitle

\section{Introduction}
Polish Sign Language (Polski J\k{e}zyk Migowy, PJM; ISO 639-3 pso) is the language of the Polish Deaf community. This dossier lists the resources that turned up in web searches. The search was brief, and several items were seen only as search-result titles, so each should be checked before use.

\section{Primary Educational Resources}
\begin{itemize}
\item ``Migaj razem z nami'': a multimedia PJM course produced for the Ministry of Education, according to the search summary. It is reported to be aimed at grades I--III and IV--VIII, to rely on illustrations and video, and to be available on the ministry site and on epodreczniki.pl. Announcement (Warsaw education authority): \url{https://www.kuratorium.waw.pl/pl/informacje/aktualnosci/15615,MIGAJ-RAZEM-Z-NAMI-multimedialny-kurs-polskiego-jezyka-migowego-PJM.pdf}
\item Ministry of Education page on a PJM course: \url{https://www.gov.pl/web/edukacja/kurs-polskiego-jezyka-migowego-pjm}
\item Press report that the ministry wants PJM taught in schools: \url{https://www.prawo.pl/oswiata/men-chce-uczyc-w-szkolach-polskiego-jezyka-migowego,135233.html}
\end{itemize}

\section{Linguistic and Reference Materials}
\begin{itemize}
\item Rutkowski and Mostowski, paper on PJM as a tool against exclusion (University of Warsaw): \url{https://www.plm.uw.edu.pl/wp-content/uploads/2012/07/Rutkowski-Mostowski-Polski-jezyk-migowy-narzedziem-przeciwdzialania-wykluczeniu.pdf}
\item Ministry of Education page on levels of proficiency in PJM: \url{https://zpe.gov.pl/a/opis/DEL51L34P}
\end{itemize}

\section{Religious and Community Resources}
\begin{itemize}
\item Polish Sign Language Council resolution 1/2024, which recommends a PJM interpreter in schools lacking PJM-competent teachers: \url{https://niepelnosprawni.gov.pl/wp-content/uploads/2025/06/Uchwala-1-2024-1743775612.docx}
\item Meeting of the minister of education with the Instytut Praw G\l uchych: \url{https://www.gov.pl/web/nauka/spotkanie-ministra-edukacji-i-nauki-z-przedstawicielami-instytutu-praw-gluchych}
\item SWPS University presentation, training in Deaf culture and PJM: \url{https://www0.swps.pl/images/DOKUMENTY/Prezentacja_-_Szkolenie_z_kultury_G%C5%82uchych_i_polskiego_j%C4%99zyka_migowego.pdf}
\end{itemize}

\section{Digital and Audio Resources}
\begin{itemize}
\item PJM A1 course programme (powiat gi\.zycki): \url{https://www.powiatgizycki.pl/wp-content/uploads/2024/11/PROGRAM-i-TERMINARZ-kursu-PJM-A1.pdf}
\end{itemize}
A corpus-based PJM dictionary was mentioned in a search summary, but no URL for it appeared in the results, so it is not listed.

\section{Curriculum Development Context}
The searches indicate that textbooks for grades 4--8 were adapted and translated into PJM in 2017--2019, but apparently were distributed to teachers rather than published openly. The ``Migaj razem z nami'' course exists for the early grades and for grades 4--8. No full K-12 subject curriculum in PJM was found. Coverage here is partial.
\end{document}